% Review draft: FF.47. Genus-zero cover rigidity, not a cofinality claim.
\subsection{The pointwise kernel fixes the genus-zero part of the site}

\begin{claim}{FF.47}{Proved}\label{clm:FFgenuszerokernel}
Let \(E/B_p\) be finite and let \(\varphi\) be a
pointwise-trivial pinned automorphism of
\(\mathbb A^1_{E,\mathrm{\acute et}}\) over \(BG_E\).
Its inverse-image functor is naturally isomorphic to the
identity on the full subcategory of quasicompact separated
etale objects whose smooth projective geometric components
have genus zero. The same assertion holds after every
finite constant extension.

This constraint is strictly stronger than peripheral and
abelianization control for an abstract free profinite
four-puncture group. Specifically, the word automorphism
\[
 x\longmapsto(xy)x(xy)^{-1},\qquad
 y\longmapsto(xy)y(xy)^{-1},\qquad z\longmapsto z
                                                        \tag{FF.47.1}
\]
of the free \emph{profinite} group on \(x,y,z\) moves
a degree-three genus-zero cover, despite fixing every
peripheral conjugacy class and abelianization. Hence it
cannot be the geometric action of such a \(\varphi\)
on a four-punctured open. This is a statement about the
full profinite word automorphism (FF.47.1); it is not an
assertion about every full automorphism whose pro-\(p\)
image equals the corresponding pro-\(p\) word automorphism.
\end{claim}

\begin{proof}
We first take a geometrically connected etale object
\(Y\to\mathbb A^1_E\) with projective completion
\(C\) of genus zero, and write
\(Y'=\varphi^*Y\). Pullback to the slice topoi gives
an equivalence \(Y'_{\mathrm{\acute et}}\to
Y_{\mathrm{\acute et}}\) over the specified base.
It induces an isomorphism of the generic Galois groups
over \(G_E\). The constant fields remain \(E\),
because the relative generic action preserves the
projection to \(G_E\). Let \(C'\) denote the smooth
projective completion of \(Y'\).

Koenigsmann's local theory identifies all geometric
inertia groups on \(C\) and \(C'\), not merely those
at the points of the two displayed opens. Killing their
closed normal span gives the etale fundamental group of
the complete curve, compatibly over \(G_E\). Hence
the geometric fundamental groups of \(C\) and \(C'\)
are isomorphic. In characteristic zero the geometric
abelianization has rank \(2g\), by
\citet[Expos\'e~XIII, Corollary~2.12]{FF_SGA1}.
It follows that \(C'\) also has genus zero.

The same local theory supplies a compatible
\(G_E\)-equivariant point bijection
\[
                         h:C'(\bar E)\longrightarrow C(\bar E).
                                                        \tag{FF.47.2}
\]
Over the displayed opens it is the actual point function
of the slice equivalence. Indeed, the specialization
functors \(S\mapsto S^{I_x}\) used in
Claim~\ref{clm:FFgenerickernel} identify the inertia
class of every actual point with precisely its transported
class. At a rational point the arithmetic conjugator can
be made geometric using the surjection of its decomposition
group onto the constant Galois group, as explained in
Claim~\ref{clm:FFallprimekernel}.

Choose a finite Galois extension \(L/E\) over which
\(C\) has a rational point \(c\). Then
\(c'=h^{-1}(c)\) is \(L\)-rational as well.
A smooth projective genus-zero curve with a rational
point is \(\mathbb P^1\). Choose coordinates on
\(C_L\) and \(C'_L\) taking \(c\) and \(c'\)
to infinity. The generic group isomorphism, together
with (FF.47.2), transports the inertia classes of the
two remaining affine lines. The site construction in
Claim~\ref{clm:FFallprimekernel}, applied after these
coordinate identifications, gives a pinned automorphism
of the affine-line etale topos whose actual point function
is \(h\) in these coordinates.

Claim~\ref{clm:FFmetabelianaffine} makes this function
affine. It therefore extends to an \(L\)-isomorphism
\(C'_L\to C_L\), realizing all of (FF.47.2).
It carries \(Y'_L\) onto \(Y_L\), because on these
opens the actual point function is the bijection of the
slice equivalence. This isomorphism descends to \(E\):
the point bijection is \(G_E\)-equivariant, so every
Galois conjugate of the \(L\)-isomorphism has the same
geometric point function and is therefore equal to it.
This is ordinary finite Galois descent of scheme maps.
We obtain a unique \(E\)-isomorphism
\[
                                  h_Y:Y'\xrightarrow{\sim}Y
                                                        \tag{FF.47.3}
\]
with the required actual point function.

It is an isomorphism over \(\mathbb A^1_E\).
The slice equivalence lies over \(\varphi\), whose
base point function is the identity. Thus both composites
from \(Y'\) to \(\mathbb A^1_E\) have the same
geometric point function, and the corresponding scheme
maps are equal. For any arrow \(Y\to Z\) between
two objects under consideration, the same pointwise
compatibility gives
\(h_Z\circ\varphi^*(Y\to Z)=(Y\to Z)\circ h_Y\).
Equality follows because the sources are reduced and
geometric points are dense. This proves naturality.

For a connected object with a finite constant field
extension, apply the same argument over that extension;
the relative structure identifies the extension on both
sides. For a disconnected object, use its finitely many
connected components. Equivalently, one can split all
components and choose their rational points over one
finite Galois constant extension, then descend the unique
pointwise isomorphisms. This proves the full first
assertion and its compatibility with further constants.

\emph{A finite quotient detecting the abstract word automorphism.}
Write \(w=(xyz)^{-1}\). In \(S_3\), assign
\[
           (x,y,z,w)\longmapsto
           \big((12),(23),(12),(13)\big).
                                                        \tag{FF.47.4}
\]
The four permutations have product one and generate
\(S_3\). By the characteristic-zero presentation of the
geometric fundamental group of a four-punctured line
in \citet[Expos\'e~XIII, Corollary~2.12]{FF_SGA1},
this gives a connected degree-three etale cover over
\(\bar E\). Each of the four branch permutations
is a transposition. Its smooth completion therefore has
\[
                         2g-2=-2\cdot3+4=-2,
\]
so has genus zero. The cover and its map descend to a
finite constant extension of \(E\), by finite
presentation.

The automorphism (FF.47.1) fixes \(xy\) and \(z\),
and therefore fixes \(w\). It conjugates the first
two permutations in (FF.47.4) by their nontrivial
three-cycle product, leaving the last two unchanged.
The original and transformed degree-three permutation
representations are not isomorphic. A simultaneous
conjugator would have to centralize both \((12)\) and
\((13)\), so would be the identity; but the first
permutation has changed. Thus (FF.47.1) moves this
genus-zero cover and contradicts the first assertion
if proposed as the action of \(\varphi\).

This finite calculation was executed in
\texttt{F\_metabelian\_checks.py}; its output verifies
the product, transitivity, Riemann--Hurwitz integer, and
the absence of any of the six possible simultaneous
conjugators. Since \(S_3\) is not a \(p\)-group,
the test concerns the stated free profinite word action
and does not identify all possible lifts of a pro-\(p\)
action. No arithmetic-compatible ghost is constructed.
\end{proof}

The genus-zero subcategory is not a generating family
for the whole curve etale site. For example, a nonempty
etale object with projective geometric genus at least
one cannot be covered by genus-zero etale curves: a
dominant map between their projective completions would
contradict Riemann--Hurwitz. Nor has any cofinality of
Galois closures of genus-zero covers been established.
The action on arbitrary positive-genus covers remains the missing
part of the birational kernel problem.
